\section{Introduction}
Hardware acceleration reduces the cost of computation within individual pipeline stages. End-to-end execution also requires converting the output of one stage into the input representation of the next, transferring that representation, and coordinating execution. These operations can account for a substantial fraction of runtime even when both stages execute on specialized hardware. Understanding their cost is therefore necessary to translate improvements in accelerator performance into improvements in application performance.

Neural--symbolic systems provide a concrete setting for studying this interface. Neural computation produces predictions from unstructured inputs, while symbolic reasoning applies explicit constraints to those predictions. In visual Sudoku, for example, a classifier recognizes handwritten givens and a SAT solver completes the grid. The interface must translate predicted digits into the solver's representation and make them available for reasoning. The Sudoku rules remain unchanged across inputs, yet a full-image host flow reconstructs and transfers the solver's data structures for every puzzle. This repeated preparation couples the cost of supplying new predictions to the size of the shared symbolic representation.

This paper investigates how much of that preparation can be removed through representation and state reuse, and what additional benefit comes from generating updates in hardware. We study an FPGA system built with ESP~\cite{esp}, combining NVDLA for int8 neural inference with SAT-Accel for conflict-driven clause-learning (CDCL) SAT solving~\cite{sataccel}. The system executes perception and reasoning on separate accelerator tiles, providing an end-to-end setting in which to measure conversion, data movement, synchronization, and solver loading.

We introduce ReasonLink, an interface organized around the separation of invariant constraints from input-dependent predictions. The shared rule formula is initialized once, and each new prediction is encoded as a solver-native fixed-decision literal. ReasonLink reads neural outputs, generates and buffers literals in hardware, and writes a compact update to the existing solver image. A complementary on-chip base reuse mechanism retains structural solver data between puzzles and restores the state modified during reasoning. Together, these mechanisms reduce repeated host preparation and internal solver loading.

Our evaluation compares \emph{CPU-generated deltas}, in which CPU software computes literals and writes the update, with \emph{hardware-generated deltas}, in which ReasonLink performs those operations. Both paths use the same fixed-decision representation and support on-chip base reuse. We compare them against full-image preparation and evaluate both sequential execution and execution that overlaps preparation of the next puzzle with the current solve. These configurations evaluate compact updates, computation placement, state reuse, and overlap.

Board measurements show that the handoff accounts for approximately 35\% of instrumented serial runtime when the host rebuilds the solver image and transfers allocated buffer regions, including unused capacity. Compact updates reduce this overhead sharply, while on-chip base reuse reduces median solver time from 10.2 to 3.1\,ms. Neural inference then dominates execution, and CPU-generated and hardware-generated deltas achieve similar end-to-end performance. The resulting characterization identifies both the repeated interface work that can be eliminated and the stage that limits further improvement on this platform.

This paper makes three contributions:
\begin{itemize}
\item An end-to-end characterization of representation conversion, data movement, and synchronization between neural and symbolic accelerators on an FPGA system.
\item ReasonLink, a hardware delta-generation interface, and an on-chip base reuse mechanism that restores solver state across inputs, validated through functional simulation, RTL simulation, and board execution.
\item A comparative evaluation of CPU-generated and hardware-generated deltas, on-chip base reuse, and execution overlap, identifying their effects on interface cost, solver time, and end-to-end performance.
\end{itemize}
